\documentclass[11pt]{article}
\usepackage{mathunal-base}
\mathunalfuente{serif}

\renewcommand{\examtitulo}{Tercer Examen Teórico de Métodos Numéricos (25\%)}
\renewcommand{\examfecha}{Mayo 14 de 2007}

\begin{document}

\examheader

\vspace{4pt}
\noindent Nombre: \dotfill\ Carné: \dotfill\ Grupo: \dotfill

\vspace{8pt}
\noindent\textbf{La interpretación del examen hace parte de la evaluación, por tanto no se admiten preguntas comenzada la prueba. Duración: 1 hora y 50 minutos.}

\vspace{10pt}
\noindent\textbf{Parte 1 Preguntas de selección múltiple. Cada una tiene un valor de 6\%.}

\noindent Escoja la única respuesta de cada pregunta y señálela haciendo una X sobre la letra correspondiente. Si en alguna pregunta señala más de una respuesta se le asigna un puntaje de cero a esa pregunta. Todas las preguntas valen lo mismo.

\begin{enumerate}
\item{} Sean $R=\{(t,y)\ /\ 0\leq t\leq 3,\ 0<y<5\}$ y $f(t,y)=2ty^2$. De la lista siguiente el menor número que le sirve a $f(t,y)$ como constante de Lipschitz con respecto a su variable $y$ en $R$ es

\vspace{4pt}
\noindent a. $15$ \hspace{1cm} b. $30$ \hspace{1cm} c. $60$ \hspace{1cm} d. $120$

\item{} Se utiliza el método de Euler para aproximar la solución del PVI $x''+6x'+9x=0,\ x(0)=4$ y $x'(0)=-4$. Si se usa $h=0.5$ y la variable vectorial es $u$, entonces el primer iterado que se obtiene por el método de Euler es $u(0.5)=$

\vspace{4pt}
\noindent
\begin{minipage}[t]{0.24\linewidth}
a. $\begin{bmatrix}2\\10\end{bmatrix}$
\end{minipage}%
\begin{minipage}[t]{0.24\linewidth}
b. $\begin{bmatrix}-10\\2\end{bmatrix}$
\end{minipage}%
\begin{minipage}[t]{0.24\linewidth}
c. $\begin{bmatrix}10\\-2\end{bmatrix}$
\end{minipage}%
\begin{minipage}[t]{0.24\linewidth}
d. $\begin{bmatrix}2\\-10\end{bmatrix}$
\end{minipage}

\item{} Utilizamos el método del disparo lineal para aproximar la solución del problema de contorno
\[
\left\{
\begin{aligned}
x'' &= 100x\\
x(0) &= 1,\ x(1)=e^{-10}
\end{aligned}
\right.
\]
\noindent Este método requiere enunciar dos PVI's. El primero de ellos requiere las siguientes condiciones iniciales:

\vspace{4pt}
\noindent a. $u(0)=0,\ u'(0)=0$ \hspace{1cm} b. $u(0)=1,\ u'(0)=0$

\noindent c. $u(0)=e^{-10},\ u'(0)=1$ \hspace{1cm} d. Ninguna de las anteriores

\item{} Para resolver el problema parabólico $u_t=5u_{xx}$ con condición inicial y condiciones de contorno conocidas, se desea utilizar un método \textbf{estable} de diferencias finitas progresivas con tamaño de paso en espacio $h=0.1$. El mayor valor que puede tomar el tamaño de paso en tiempo $k$ es:

\vspace{4pt}
\noindent a. $0.1$. \hspace{1cm} b. $0.01$. \hspace{1cm} c. $0.001$. \hspace{1cm} d. $0.0001$.

\item{} Se desea resolver el problema $u_t=5u_{xx}$ con la condición inicial $u(x,0)=\sin(\pi x)$ para $0\leq x\leq 1$ y las condiciones de borde $u(0,t)=u(1,t)=0$ para $t>0$. Se utiliza un método de diferencias finitas progresivas con tamaño de paso en espacio $h=\frac{1}{8}=0.125$ y tamaño de paso en tiempo $k=\frac{1}{800}=0.00125$. La aproximación que se obtiene para $u$ evaluada en $x=\frac{1}{4}$ y $t=0.00125$ es:

\vspace{4pt}
\noindent
\begin{minipage}[t]{0.48\linewidth}
a. $0.2\sin\left(\frac{\pi}{4}\right)+0.4\left(\sin\left(\frac{\pi}{8}\right)+\sin\left(\frac{3\pi}{8}\right)\right)$
\end{minipage}%
\begin{minipage}[t]{0.48\linewidth}
b. $0.2\sin\left(\frac{\pi}{8}\right)+0.4\sin\left(\frac{\pi}{4}\right)$
\end{minipage}

\vspace{6pt}
\noindent
\begin{minipage}[t]{0.48\linewidth}
c. $0.4\sin\left(\frac{\pi}{4}\right)+0.2\left(\sin\left(\frac{\pi}{8}\right)+\sin\left(\frac{3\pi}{8}\right)\right)$
\end{minipage}%
\begin{minipage}[t]{0.48\linewidth}
d. Ninguna de las anteriores.
\end{minipage}
\end{enumerate}

\vspace{10pt}
\noindent\textbf{Parte 2 Preguntas abiertas}

\begin{enumerate}
\setcounter{enumi}{5}
\item{} Consideremos el problema de valor inicial $y'=e^{-2t}-2y,\ y(0)=1$.
\begin{enumerate}
\item{} (10\%) Demuestre que este problema tiene una única solución.

\vspace{3cm}

\item{} (15\%) ¿Cuál es el valor que se obtiene con el método de Heun como aproximación para $y(0.4)$ utilizando $h=0.2$?

\vspace{4cm}
\end{enumerate}

\item{} (25\%) Aplique el método de Runge Kutta 4 con $h=0.4$ para aproximar el valor de $u_3(0.4)$ para el sistema
\[
\left\{
\begin{aligned}
u_1' &= u_2-u_3+t & 0\leq t\leq 1\\
u_2' &= 3t^2 & 0\leq t\leq 1\\
u_3' &= u_2+e^{-t} & 0\leq t\leq 1
\end{aligned}
\right.
\qquad u_1(0)=1,\quad u_2(0)=1,\quad u_3(0)=-1
\]

\vspace{6cm}

\item{} Considere el problema de valores en la frontera o de contorno
\[
\left\{
\begin{aligned}
x'' &= -3x'+2x+2t+3 & 0\leq t\leq 1\\
x(0) &= 2 & x(1)=1
\end{aligned}
\right.
\]
\begin{enumerate}
\item{} (10\%) Demuestre que tiene solución única.

\vspace{4cm}

\item{} (10\%) Enuncie los dos problemas de valor inicial que se deben resolver para utilizar el método del disparo lineal.

\vspace{4cm}
\end{enumerate}
\end{enumerate}

\examfooter
\end{document}
